% Copyright (c) 2026 Eugene Joseph M. Ragasa.
% SPDX-License-Identifier: Apache-2.0
% Author: Eugene Joseph M. Ragasa
% Formal proof completed: 2026-10-19
\documentclass[11pt]{article}
\usepackage[T1]{fontenc}\usepackage{amsmath,amssymb}\usepackage[margin=1in]{geometry}\usepackage[hidelinks]{hyperref}
\title{Proof Dossier: Exact Rational Enclosure of $1/\sqrt{2}$}\author{Eugene Joseph M. Ragasa}\date{Formal proof completed: 2026-10-19}
\begin{document}\maketitle
\section{Identity and status}
ID \texttt{BG-CERT-THM-009}; status \texttt{proved}. Lean declaration:
\par\noindent{\tiny\path{ProjectKoios.BandGauge.Certificates.}}\\
{\tiny\path{sevenTenths_lt_inv_sqrtTwo_lt_seventyOneHundredths}}.
\section{Exact statement}
The kernel checks
\[
 \frac{7}{10}<\frac{1}{\sqrt{2}}<\frac{71}{100}.
\]
\section{Proof and boundary}
Positivity of $\sqrt2$ permits positive-denominator comparison. Squaring the resulting
nonnegative sides and using $(\sqrt2)^2=2$ reduces both inequalities to exact rational
arithmetic. This encloses the algebraic center appearing near Paper A manuscript line 857
at owner commit
\begin{center}\small\texttt{99a46cb59d0c6f27ada7d709f194996ef9619d9f}.\end{center}
It does not prove that a concrete field has that center. The kernel axiom report is:
\begin{itemize}
\item \texttt{propext};
\item \texttt{Classical.choice}; and
\item \texttt{Quot.sound}.
\end{itemize}
There are no admissions. The proof is AI-assisted and kernel-checked.
\end{document}
